\documentclass[11pt]{article}
\usepackage[utf8]{inputenc}
\usepackage{amsmath,amssymb,booktabs,geometry,verbatim}
\usepackage{CJKutf8}
\geometry{margin=1in}
\usepackage[hidelinks]{hyperref}

\title{CSP-J 2020 (Junior)}
\date{2020-11-07, 8:30--12:00 (3.5 hours)}

\begin{document}
\maketitle

\noindent
This file is tailored to QOJ: every problem reads from standard input and
writes to standard output (the original paper used the files
\texttt{<code>.in} / \texttt{<code>.out} named after each problem).

\noindent
Problems of the second round of the 2020 CCF Certified Software Professional
(Junior group), in contest order.
Every problem has time limit 1 s and memory limit 256 MB (as shown on QOJ),
and is worth 100 points.

\clearpage
\section*{Problem 1: Excellent Split (\begin{CJK}{UTF8}{gbsn}优秀的拆分\end{CJK})
  \hfill{\normalsize\href{https://qoj.ac/problem/18071}{QOJ 18071}}}

\noindent
Time limit: 1 s.
Memory limit: 256 MB.

\subsection*{Problem description}

In general, a positive integer can be split into a sum of several positive
integers.
For example, $1 = 1$, $10 = 1 + 2 + 3 + 4$, and so on.

A particular split of a positive integer $n$ is called \emph{excellent}
if and only if, in this split, $n$ is decomposed into several distinct
positive-integer powers of $2$.
Note that a number $x$ can be expressed as a positive-integer power of $2$
if and only if $x$ can be obtained by multiplying a positive number of $2$s
together.

For example, $10 = 8 + 2 = 2^3 + 2^1$ is an excellent split.
However, $7 = 4 + 2 + 1 = 2^2 + 2^1 + 2^0$ is not an excellent split,
because $1$ is not a positive-integer power of $2$.

Now, given a positive integer $n$, you need to determine whether, among all
splits of this number, there exists an excellent split.
If one exists, please give the concrete split.

\subsection*{Input}

The input file has only one line, containing a positive integer $n$,
the number to be examined.

\subsection*{Output}

If, among all splits of this number, there exists an excellent split, then
you need to output every number in this split from largest to smallest,
with adjacent numbers separated by a single space.
It can be proved that, once the order of the numbers in the split is fixed,
the split is unique.

If no excellent split exists, output ``\texttt{-1}'' (without the quotation
marks).

\subsection*{Sample 1 input}

\begin{verbatim}
6
\end{verbatim}

\subsection*{Sample 1 output}

\begin{verbatim}
4 2
\end{verbatim}

\subsection*{Explanation of sample 1}

$6 = 4 + 2 = 2^2 + 2^1$ is an excellent split.
Note that $6 = 2 + 2 + 2$ is not an excellent split, because the $3$
numbers in the split do not satisfy the requirement that all numbers be
pairwise distinct.

\subsection*{Sample 2 input}

\begin{verbatim}
7
\end{verbatim}

\subsection*{Sample 2 output}

\begin{verbatim}
-1
\end{verbatim}

\subsection*{Sample 3}

Large sample, not printed here;
provided as the downloadable files \texttt{power/power3.in} and \texttt{power/power3.ans}.

\subsection*{Constraints}

For $20\%$ of the data, $n \le 10$.

For another $20\%$ of the data, $n$ is guaranteed to be odd.

For another $20\%$ of the data, $n$ is guaranteed to be a positive-integer
power of $2$.

For $80\%$ of the data, $n \le 1024$.

For $100\%$ of the data, $1 \le n \le 1 \times 10^7$.

\clearpage
\section*{Problem 2: Live Awards (\begin{CJK}{UTF8}{gbsn}直播获奖\end{CJK})
  \hfill{\normalsize\href{https://qoj.ac/problem/18072}{QOJ 18072}}}

\noindent
Time limit: 1 s.
Memory limit: 256 MB.

\subsection*{Problem description}

NOI2130 is about to be held.
To make it more entertaining, CCF has decided to announce the contestants'
scores one by one and to broadcast the real-time award cutoff score live.
The award rate of this competition is $w\%$, that is, the lowest score among
the contestants currently ranked in the top $w\%$ is the real-time cutoff
score.

More specifically, if the scores of $p$ contestants have been announced so
far, then the current planned number of award winners is
$\max(1, \lfloor p \times w\% \rfloor)$,
where $w$ is the award percentage, $\lfloor x \rfloor$ denotes $x$ rounded
down to an integer, and $\max(x, y)$ denotes the larger of $x$ and $y$.
If some contestants have equal scores, then all contestants tied at that
score receive an award, so the actual number of award winners may be larger
than planned.

As a technical staff member of the judging team, please help CCF write a
live-broadcast program.

\subsection*{Input}

Line $1$ contains two positive integers $n, w$, the total number of
contestants and the award rate, respectively.

Line $2$ contains $n$ non-negative integers, the contestants' scores in
the order in which they are announced one by one.

\subsection*{Output}

Only one line, containing $n$ non-negative integers, which are, in order,
the real-time award cutoff scores after each contestant's score is
announced one by one.
Adjacent integers are separated by a single space.

\subsection*{Sample 1 input}

\begin{verbatim}
10 60
200 300 400 500 600 600 0 300 200 100
\end{verbatim}

\subsection*{Sample 1 output}

\begin{verbatim}
200 300 400 400 400 500 400 400 300 300
\end{verbatim}

\subsection*{Explanation of sample 1}

\begin{center}
\begin{tabular}{|l|c|c|c|c|c|c|c|c|c|c|}
\hline
Number of contestants judged & 1 & 2 & 3 & 4 & 5 & 6 & 7 & 8 & 9 & 10 \\
\hline
Planned number of award winners & 1 & 1 & 1 & 2 & 3 & 3 & 4 & 4 & 5 & 6 \\
\hline
\end{tabular}
\end{center}

\begin{center}
\begin{tabular}{|l|c|c|c|c|c|c|c|c|c|c|}
\hline
Scores of the judged contestants,
  & \textbf{200} & \textbf{300} & \textbf{400} & 500 & 600 & 600
  & 600 & 600 & 600 & 600 \\
sorted from high to low
  &  & 200 & 300 & \textbf{400} & 500 & 600 & 600 & 600 & 600 & 600 \\
(the cutoff score is marked
  &  &  & 200 & 300 & \textbf{400} & 500 & 500 & 500 & 500 & 500 \\
in \textbf{bold})
  &  &  &  & 200 & 300 & \textbf{400} & \textbf{400} & \textbf{400}
  & 400 & 400 \\
  &  &  &  &  & 200 & 300 & 300 & 300 & \textbf{300} & 300 \\
  &  &  &  &  &  & 200 & 200 & 300 & 300 & \textbf{300} \\
  &  &  &  &  &  &  & 0 & 200 & 200 & 200 \\
  &  &  &  &  &  &  &  & 0 & 200 & 200 \\
  &  &  &  &  &  &  &  &  & 0 & 100 \\
  &  &  &  &  &  &  &  &  &  & 0 \\
\hline
\end{tabular}
\end{center}

Note that after the score of the 9th contestant is announced, the planned
number of award winners is $5$, but because of ties, $6$ contestants
actually receive awards.

\subsection*{Sample 2 input}

\begin{verbatim}
10 30
100 100 600 100 100 100 100 100 100 100
\end{verbatim}

\subsection*{Sample 2 output}

\begin{verbatim}
100 100 600 600 600 600 100 100 100 100
\end{verbatim}

\subsection*{Sample 3}

Large sample, not printed here;
provided as the downloadable files \texttt{live/live3.in} and \texttt{live/live3.ans}.

\subsection*{Constraints}

\begin{center}
\begin{tabular}{|c|l|}
\hline
Test case number & $n$ \\
\hline
1$\sim$3 & $= 10$ \\
\hline
4$\sim$6 & $= 500$ \\
\hline
7$\sim$10 & $= 2000$ \\
\hline
11$\sim$17 & $= 10000$ \\
\hline
18$\sim$20 & $= 100000$ \\
\hline
\end{tabular}
\end{center}

For all test cases, every contestant's score is a non-negative integer not
exceeding $600$, and the award percentage $w$ is a positive integer with
$1 \le w \le 99$.

When computing the planned number of award winners, if a floating-point
variable (such as \texttt{float} or \texttt{double} in C/C++, or
\texttt{real}, \texttt{double}, \texttt{extended}, etc.\ in Pascal) is used
to store the award ratio $w\%$, then the result of computing
$5 \times 60\%$ may be $3.000001$ or may be $2.999999$, and the result
after rounding down is not determined.
Therefore, it is recommended to use only integer variables, so as to
compute the exact value.

\clearpage
\section*{Problem 3: Expression (\begin{CJK}{UTF8}{gbsn}表达式\end{CJK})
  \hfill{\normalsize\href{https://qoj.ac/problem/18073}{QOJ 18073}}}

\noindent
Time limit: 1 s.
Memory limit: 256 MB.

\subsection*{Problem description}

Little C is keen on studying mathematical logic.
One day, he discovered a special kind of logical expression.
In this kind of logical expression, all operands are variables, their values
can only be $0$ or $1$, and the operations are carried out from left to
right.
If the expression contains parentheses, the value of the sub-expression
inside the parentheses is computed first.
In particular, this kind of expression has exactly the following
operations:
\begin{enumerate}
\item AND: $a$ \texttt{\&} $b$.
  The value of this expression is $1$ if and only if the values of both
  $a$ and $b$ are $1$.
  In all other cases the value of this expression is $0$.
\item OR: $a$ \texttt{|} $b$.
  The value of this expression is $0$ if and only if the values of both
  $a$ and $b$ are $0$.
  In all other cases the value of this expression is $1$.
\item NOT: \texttt{!} $a$.
  The value of this expression is $1$ if and only if the value of $a$ is
  $0$.
  In all other cases the value of this expression is $0$.
\end{enumerate}

Little C wants to know: given a logical expression and the initial value of
each of its operands, what is the value of the original expression after
the value of one particular operand is negated?

To simplify the processing of the expression, we make the following
conventions.

The expression is given in postfix notation.
Postfix notation is defined as follows:
\begin{enumerate}
\item If $E$ is an operand, then the postfix expression of $E$ is $E$
  itself.
\item If $E$ is an expression of the form $E_1 \; op \; E_2$, where $op$
  is any binary operator whose precedence is not higher than that of the
  operators outside parentheses in $E_1$ and $E_2$, then the postfix form
  of $E$ is $E_1' \; E_2' \; op$,
  where $E_1'$ and $E_2'$ are the postfix forms of $E_1$ and $E_2$,
  respectively.
\item If $E$ is an expression of the form $(E_1)$, then the postfix form of
  $E_1$ is the postfix form of $E$.
\end{enumerate}

Also, for convenience, in the input:
\begin{enumerate}
\item[a)] There is one space on each side of the AND operator (\texttt{\&}),
  the OR operator (\texttt{|}) and the NOT operator (\texttt{!}), but
  there is no space at the end of the expression.
\item[b)] An operand consists of the lowercase letter \texttt{x} followed
  by a positive integer, the positive integer being the index of the
  variable.
  For example, \texttt{x10} denotes the variable $x_{10}$ with index $10$.
  It is guaranteed that every variable appears exactly once in the
  expression.
\end{enumerate}

\subsection*{Input}

The first line contains a string $s$, the expression described above.

The second line contains a positive integer $n$, the number of variables in
the expression.
The indices of the variables in the expression are $1, 2, \ldots, n$.

The third line contains $n$ integers, the $i$-th of which is the initial
value of the variable $x_i$.

The fourth line contains a positive integer $q$, the number of queries.

Each of the next $q$ lines contains a positive integer, the index of the
variable to be negated.
Note that the modification in each query is temporary, that is, the
modification in an earlier query has no effect on later queries.

It is guaranteed that the input expression is valid.
The initial value of each variable is $0$ or $1$.

\subsection*{Output}

The output consists of $q$ lines, each containing a $0$ or a $1$, the value
of the expression under the corresponding query.

\subsection*{Sample 1 input}

\begin{verbatim}
x1 x2 & x3 |
3
101
3
1
2
3
\end{verbatim}

\subsection*{Sample 1 output}

\begin{verbatim}
1
1
0
\end{verbatim}

\subsection*{Explanation of sample 1}

The infix form of this postfix expression is
$(x_1 \;\texttt{\&}\; x_2) \;\texttt{|}\; x_3$.

For the first query, the value of $x_1$ is negated.
Then the values of the three operands are, in order, $0$, $0$, $1$.
The value of the original expression is
$(0 \;\texttt{\&}\; 0) \;\texttt{|}\; 1 = 1$.

For the second query, the value of $x_2$ is negated.
Then the values of the three operands are, in order, $1$, $1$, $1$.
The value of the original expression is
$(1 \;\texttt{\&}\; 1) \;\texttt{|}\; 1 = 1$.

For the third query, the value of $x_3$ is negated.
Then the values of the three operands are, in order, $1$, $0$, $0$.
The value of the original expression is
$(1 \;\texttt{\&}\; 0) \;\texttt{|}\; 0 = 0$.

\subsection*{Sample 2 input}

\begin{verbatim}
x1 ! x2 x4 | x3 x5 ! & & ! &
5
01011
3
1
3
5
\end{verbatim}

\subsection*{Sample 2 output}

\begin{verbatim}
0
1
1
\end{verbatim}

\subsection*{Explanation of sample 2}

The infix form of this expression is
$(\texttt{!}\, x_1) \;\texttt{\&}\;
 (\texttt{!}\, ((x_2 \;\texttt{|}\; x_4) \;\texttt{\&}\;
 (x_3 \;\texttt{\&}\; (\texttt{!}\, x_5))))$.

\subsection*{Sample 3}

Large sample, not printed here;
provided as the downloadable files \texttt{expr/expr3.in} and \texttt{expr/expr3.ans}.

\subsection*{Constraints}

For $20\%$ of the data, the expression contains only AND operations
(\texttt{\&}) or only OR operations (\texttt{|}).

For another $30\%$ of the data, $|s| \le 1000$, $q \le 1000$,
$n \le 1000$.

For another $20\%$ of the data, the initial values of the variables are all
$0$ or all $1$.

For $100\%$ of the data, $1 \le |s| \le 1 \times 10^6$,
$1 \le q \le 1 \times 10^5$, $2 \le n \le 1 \times 10^5$.

Here $|s|$ denotes the length of the string $s$.

\clearpage
\section*{Problem 4: Taking Numbers on a Grid (\begin{CJK}{UTF8}{gbsn}方格取数\end{CJK})
  \hfill{\normalsize\href{https://qoj.ac/problem/18074}{QOJ 18074}}}

\noindent
Time limit: 1 s.
Memory limit: 256 MB.

\subsection*{Problem description}

There is an $n \times m$ grid, and every cell contains an integer.
A little bear wants to walk from the top-left corner of the grid to the
bottom-right corner.
In each step it may move only one cell up, down or right, it may not pass
through a cell it has already visited, and it may not leave the grid.
The bear takes the integers in all cells it passes through.
Find the maximum possible sum of the integers it takes.

\subsection*{Input}

Line $1$ contains two positive integers $n, m$.

Each of the next $n$ lines contains $m$ integers, the integers in the
cells, in order.

\subsection*{Output}

One integer, the maximum possible sum of the integers the bear can take.

\subsection*{Sample 1 input}

\begin{verbatim}
3 4
1 -1 3 2
2 -1 4 -1
-2 2 -3 -1
\end{verbatim}

\subsection*{Sample 1 output}

\begin{verbatim}
9
\end{verbatim}

\subsection*{Explanation of sample 1}

In the figures below, the cells on the path are shown in bold and the
arrows show the direction of each step.

\begin{center}
\begin{tabular}{ccccccc}
\textbf{1} & & $-1$ & & \textbf{3} & $\rightarrow$ & \textbf{2} \\
$\downarrow$ & & & & $\uparrow$ & & $\downarrow$ \\
\textbf{2} & $\rightarrow$ & $\mathbf{-1}$ & $\rightarrow$ & \textbf{4}
  & & $\mathbf{-1}$ \\
 & & & & & & $\downarrow$ \\
$-2$ & & 2 & & $-3$ & & $\mathbf{-1}$ \\
\end{tabular}
\end{center}

With the walk shown above, the sum of the numbers taken is
$1 + 2 + (-1) + 4 + 3 + 2 + (-1) + (-1) = 9$,
and it can be proved that this is the maximum.

\begin{center}
\begin{tabular}{ccccccc}
\textbf{1} & & $-1$ & & \textbf{3} & $\rightarrow$ & \textbf{2} \\
$\downarrow$ & & & & $\uparrow$ & & $\downarrow$ \\
\textbf{2} & $\rightarrow$ & $\mathbf{-1}$ & $\rightarrow$ & \textbf{4}
  & & $\mathbf{-1}$ \\
 & & $\downarrow\uparrow$ & & & & $\downarrow$ \\
$-2$ & & \textbf{2} & & $-3$ & & $\mathbf{-1}$ \\
\end{tabular}
\end{center}

Note that the walk shown above is invalid, because the cell in row $2$,
column $2$ is visited twice, and according to the problem statement a cell
that has already been visited may not be passed through again.

\begin{center}
\begin{tabular}{ccccccc}
\textbf{1} & & $-1$ & & \textbf{3} & $\rightarrow$ & \textbf{2} \\
$\downarrow$ & & & & $\uparrow$ & & \\
\textbf{2} & $\rightarrow$ & $\mathbf{-1}$ & $\rightarrow$ & \textbf{4}
  & & $-1$ \\
 & & & & & & \\
$-2$ & & 2 & & $-3$ & & $-1$ \\
\end{tabular}
\end{center}

Moreover, the walk shown above is also invalid, because it does not reach
the end point at the bottom-right corner.

\subsection*{Sample 2 input}

\begin{verbatim}
2 5
-1 -1 -3 -2 -7
-2 -1 -4 -1 -2
\end{verbatim}

\subsection*{Sample 2 output}

\begin{verbatim}
-10
\end{verbatim}

\subsection*{Explanation of sample 2}

\begin{center}
\begin{tabular}{ccccccccc}
$\mathbf{-1}$ & $\rightarrow$ & $\mathbf{-1}$ & $\rightarrow$
  & $\mathbf{-3}$ & $\rightarrow$ & $\mathbf{-2}$ & & $-7$ \\
 & & & & & & $\downarrow$ & & \\
$-2$ & & $-1$ & & $-4$ & & $\mathbf{-1}$ & $\rightarrow$
  & $\mathbf{-2}$ \\
\end{tabular}
\end{center}

With the walk shown above, the sum of the numbers taken is
$(-1) + (-1) + (-3) + (-2) + (-1) + (-2) = -10$,
and it can be proved that this is the maximum.
Therefore, please note that the maximum sum of the numbers taken may also
be negative.

\subsection*{Sample 3}

Large sample, not printed here;
provided as the downloadable files \texttt{number/number3.in} and \texttt{number/number3.ans}.

\subsection*{Constraints}

For $20\%$ of the data, $n, m \le 5$.

For $40\%$ of the data, $n, m \le 50$.

For $70\%$ of the data, $n, m \le 300$.

For $100\%$ of the data, $1 \le n, m \le 1000$.
The absolute value of every integer in the grid does not exceed $10^4$.

\section*{Source}

The statements were taken from the official CSP-J 2020 problem PDFs served
by QOJ (problem pages, ``Statement'' tab):
\begin{itemize}
\item Problem 1: \texttt{https://qoj.ac/problem/18071}
\item Problem 2: \texttt{https://qoj.ac/problem/18072}
\item Problem 3: \texttt{https://qoj.ac/problem/18073}
\item Problem 4: \texttt{https://qoj.ac/problem/18074}
\end{itemize}
Contest page: \texttt{https://qoj.ac/contest/2860} (CSP-J 2020).

\end{document}
